\section{Problem Formulation}
\label{sec:theoretical_problem_formulation}

The study is formulated within an episodic reinforcement-learning setting.
An environment is represented as a Markov decision process with observation space $\mathcal{O}$, action space $\mathcal{A}$, transition distribution $P(o_{t+1}\mid o_t,a_t)$, extrinsic reward function $r^{\mathrm{ext}}(o_t,a_t,o_{t+1})$, discount factor $\gamma\in[0,1)$, and finite episode horizon $H$.
At each time step $t$, the agent receives an observation $o_t\in\mathcal{O}$, samples an action $a_t\sim\pi_\theta(\cdot\mid o_t)$ from a stochastic policy, receives an extrinsic reward $r_t^{\mathrm{ext}}$, and transitions to $o_{t+1}$.
The notation uses observations rather than privileged environment states because the exploration mechanism is defined only from information available to the agent.

The primary task objective is the maximization of expected discounted extrinsic return,
\begin{equation}
    J_{\mathrm{ext}}(\theta)
    = \mathbb{E}_{\tau\sim\pi_\theta}
    \left[\sum_{t=0}^{H-1}\gamma^t r_t^{\mathrm{ext}}\right],
\end{equation}
where $\tau=(o_0,a_0,r_0^{\mathrm{ext}},o_1,\ldots)$ denotes a trajectory induced by the policy and environment.
This objective defines the performance criterion of the task.
Intrinsic rewards are introduced only as an auxiliary training signal for exploration and do not replace extrinsic return as the quantity used to evaluate learned behavior.

Sparse-reward and delayed-reward tasks make direct optimization of $J_{\mathrm{ext}}$ difficult because many early trajectories contain little information about useful behavior.
The theoretical role of intrinsic motivation is to modify the training signal so that the policy visits informative parts of the observation--action space before extrinsic reward becomes frequent.
The resulting optimization process is therefore governed by a shaped reward during training, while the success of the learned policy is judged by the original environment reward.
This distinction is important because arbitrary reward shaping does not generally preserve optimal policies, unlike restricted potential-based shaping forms \parencite{ng-1999}.

The central problem addressed by the framework is to define an intrinsic reward $r_t^{\mathrm{int}}$ that rewards productive exploration rather than persistent surprise.
A productive transition is characterized by two properties.
First, it should produce a meaningful change in a learned feature representation, indicating that the action affected the agent's represented state.
Second, it should occur in a region where the predictive model is improving, indicating reducible uncertainty rather than irreducible noise.
The following sections formalize these properties through latent dynamics, feature-space impact, region-local learning progress, and a gate that suppresses intrinsically rewarding but unproductive regions.